\chapter{Introduction to C}

You already know how to write programs in Python.  In the next few
chapters, you are going to learn C.  Your Python knowledge will make
this easier, because C has the same fundamental ideas: variables,
functions, loops, and conditionals.  The syntax looks different, and
there are a few genuinely new concepts, but you will recognize the shape
of everything you write.

Why learn C?  Two reasons.  First, C is used wherever performance
matters: operating systems, game engines, embedded systems, and
scientific computing.  Second, and more important for this course, C
forces you to think about memory explicitly.  Python quietly manages
memory on your behalf; C makes you do it yourself.  Once you understand
how C uses memory, you will also understand what Python is doing
underneath, and you will be ready to implement data structures from
scratch.

\section{Compiled vs.\ Interpreted}

When you run a Python program, the Python interpreter reads your
\texttt{.py} file and executes it line by line.  There is no separate
step between writing and running.

C works differently.  Before a C program can run, a
\emph{compiler}\index{compiler} must translate it into machine code,
the raw instructions that your CPU executes directly.  The compiler
reads your \texttt{.c} file, checks it for errors, and produces an
\emph{executable} file.  You then run that executable.

\begin{figure}[h]
\centering
\begin{tabular}{ccccc}
\texttt{hello.c} & $\xrightarrow{\text{compiler}}$ &
\texttt{hello} (executable) & $\xrightarrow{\text{run}}$ & output
\end{tabular}
\caption{The compile-then-run cycle}
\end{figure}

Because the compiler sees your entire program before anything runs, it
can catch many mistakes, like using a variable before declaring it or
passing the wrong type to a function, that Python would only discover
at runtime.

To compile a C file called \texttt{hello.c}, you run:

\begin{minted}{bash}
gcc hello.c -o hello
\end{minted}

The \texttt{-o hello} flag names the output executable \texttt{hello}.
Then you run it:

\begin{minted}{bash}
./hello
\end{minted}

\section{Your First C Program}

Create a file called \texttt{hello.c} and type the following:

\begin{minted}{c}
#include <stdio.h>

int main()
{
    printf("Hello, world!\n");
    return 0;
}
\end{minted}

Compile and run it.  You should see:

\begin{minted}{c}
Hello, world!
\end{minted}

Let's look at each piece.

\texttt{\#include <stdio.h>} tells the compiler to include the standard
input/output library, which provides \texttt{printf} for printing.
Think of it like Python's \texttt{import} statement.

\texttt{int main()} is the function where your program starts.  Every
C program must have exactly one \texttt{main} function.  The \texttt{int}
means it returns an integer; \texttt{return 0} signals to the operating
system that the program finished successfully.

\texttt{printf("Hello, world!\textbackslash n");} prints a line of text.
The \texttt{\textbackslash n} is the newline character.

Notice the semicolons at the end of statements, and the curly braces
\texttt{\{ \}} around the body of \texttt{main}.  In Python, indentation
defines blocks; in C, curly braces do.

\section{Types and Variables}

In Python, you write:

\begin{minted}{python}
x = 5
y = 3.14
\end{minted}

Python figures out the type of each variable automatically.

In C, you must declare the type of every variable when you create it:

\begin{minted}{c}
int x = 5;
double y = 3.14;
\end{minted}

This is called \emph{static typing}\index{static typing}.  Once you
declare a variable as \texttt{int}, it can only hold integers.  The
compiler enforces this and will refuse to compile code that mixes types
incorrectly.

The common C types are:\index{C types}

\begin{description}
\item[\texttt{int}] A whole number: \texttt{0}, \texttt{42},
  \texttt{-7}.  Typically 4 bytes.
\item[\texttt{double}] A floating-point number: \texttt{3.14},
  \texttt{-0.001}.  Typically 8 bytes.  (There is also \texttt{float},
  which uses 4 bytes and has less precision.)
\item[\texttt{char}] A single character: \texttt{'A'}, \texttt{'z'},
  \texttt{'7'}.  Use single quotes.
\item[\texttt{int}] Booleans in C are just integers: \texttt{0} means
  false, and any non-zero value means true.  You can include
  \texttt{<stdbool.h>} to get the names \texttt{true} and \texttt{false}.
\end{description}

You can also declare a variable without immediately giving it a value:

\begin{minted}{c}
int count;      /* declared but not initialized */
count = 10;     /* now it has a value */
\end{minted}

Warning: an uninitialized variable in C contains whatever random bytes
happen to be in that memory location.  Unlike Python, C will not raise
an error if you read an uninitialized variable; it will just give you
garbage.  Always initialize your variables.

Note that C uses \texttt{/* ... */} for comments, not \texttt{\#}.
You can also use \texttt{//} for single-line comments in modern C.

\section{Printing with \texttt{printf}}

\texttt{printf} uses format specifiers to print different types:

\begin{minted}{c}
int age = 17;
double height = 1.75;
printf("Age: %d, Height: %f\n", age, height);
\end{minted}

Output:
\begin{minted}{c}
Age: 17, Height: 1.750000
\end{minted}

Common format specifiers:

\begin{center}
\begin{tabular}{ll}
\hline
\textbf{Specifier} & \textbf{Type} \\
\hline
\texttt{\%d}  & \texttt{int} \\
\texttt{\%f}  & \texttt{double} or \texttt{float} \\
\texttt{\%c}  & \texttt{char} \\
\texttt{\%s}  & string (char pointer) \\
\hline
\end{tabular}
\end{center}

\section{Conditionals and Loops}

The logic is the same as Python; only the syntax changes.

\subsection{if / else}

\begin{minted}{c}
int score = 85;

if (score >= 90) {
    printf("A\n");
} else if (score >= 80) {
    printf("B\n");
} else {
    printf("C or below\n");
}
\end{minted}

The condition must be in parentheses \texttt{( )}, and each branch is
surrounded by curly braces \texttt{\{ \}}.

\subsection{for and while loops}

\begin{minted}{c}
for (int i = 0; i < 5; i++) {
    printf("%d\n", i);
}
\end{minted}

A C \texttt{for} loop has three parts separated by semicolons: the
initializer (\texttt{int i = 0}), the condition (\texttt{i < 5}), and
the update (\texttt{i++}).  The \texttt{++} increments by 1, the same
as \texttt{i = i + 1}.

\begin{minted}{c}
int i = 0;
while (i < 5) {
    printf("%d\n", i);
    i++;
}
\end{minted}

\section{Functions}

In C, you must declare the return type and the type of each parameter:

\begin{minted}{c}
double square(double x)
{
    return x * x;
}

int main()
{
    printf("%f\n", square(4.0));   /* prints 16.000000 */
    printf("%f\n", square(2.5));   /* prints 6.250000  */
    return 0;
}
\end{minted}

If a function returns nothing, its return type is \texttt{void}:

\begin{minted}{c}
void printGreeting(char *name)
{
    printf("Hello, %s!\n", name);
}
\end{minted}

Functions must be defined \emph{before} they are called, or you must
provide a \emph{forward declaration} (just the signature, ending in a
semicolon) above \texttt{main}:

\begin{minted}{c}
double square(double x);   /* forward declaration */

int main()
{
    printf("%f\n", square(3.0));
    return 0;
}

double square(double x)    /* full definition can come later */
{
    return x * x;
}
\end{minted}

\section{The Stack: Blackboards for Functions}
\index{stack}\index{stack frame}

Every function, while it is running, needs somewhere to store its
local variables.  C gives each function a \emph{frame}: a reserved
chunk of memory that holds all of that function's local variables and
parameters.

Think of a frame as a \textbf{blackboard}.  When a function starts
running, it gets a fresh blackboard.  It can write anything it wants
on that blackboard, create variables, change their values.  When the
function finishes, the blackboard is \emph{erased and discarded}.
Local variables do not survive after their function returns.

Where do these frames live?  In a region of memory called the
\emph{stack}\index{stack}.  The stack works exactly like a physical
stack of objects: the most recently added item is always on top.  When
\texttt{main} calls a function, that function's frame is pushed onto
the top of the stack.  When the function returns, its frame is popped
off, and \texttt{main}'s frame is on top again.

Here is a concrete example:

\begin{minted}{c}
#include <stdio.h>

int cookingMinutes(int weightPounds)
{
    int minutes = 15 + 15 * weightPounds;
    return minutes;
}

int main()
{
    int totalWeight = 10;
    int gibletsWeight = 1;
    int turkeyWeight = totalWeight - gibletsWeight;
    int result = cookingMinutes(turkeyWeight);
    printf("Cook for %d minutes.\n", result);
    return 0;
}
\end{minted}

When \texttt{main} calls \texttt{cookingMinutes(9)}, the stack looks
like this:

\begin{center}
\begin{tabular}{r|l|}
\hline
\textbf{cookingMinutes()} & \texttt{weightPounds = 9} \\
                          & \texttt{minutes = 150} \\
\hline
\textbf{main()}           & \texttt{totalWeight = 10} \\
                          & \texttt{gibletsWeight = 1} \\
                          & \texttt{turkeyWeight = 9} \\
                          & \texttt{result = ?} \\
\hline
\end{tabular}
\end{center}

\texttt{main}'s blackboard sits below; \texttt{cookingMinutes}'s
blackboard sits on top of it.  The variable \texttt{minutes} inside
\texttt{cookingMinutes} is completely separate from anything in
\texttt{main}.  Each function only sees its own blackboard.

When \texttt{cookingMinutes} returns 150, its frame is discarded.
\texttt{main}'s frame resumes, and \texttt{result} is filled in with 150.

\subsection{Parameters are copies}

When you call \texttt{cookingMinutes(turkeyWeight)}, C copies the
value of \texttt{turkeyWeight} into the new parameter
\texttt{weightPounds} on \texttt{cookingMinutes}'s blackboard.
Changing \texttt{weightPounds} inside the function would not change
\texttt{turkeyWeight} in \texttt{main}.  They are separate variables.
This is called \emph{pass-by-value}\index{pass-by-value}.

\section{Addresses and Pointers}
\index{pointer}

Every variable in your program lives at some address in memory.  You
can find out the address of a variable using the \texttt{\&} operator,
called the \emph{address-of} operator\index{address-of operator}:

\begin{minted}{c}
#include <stdio.h>

int main()
{
    int i = 17;
    printf("i stores its value at %p\n", (void *)&i);
    return 0;
}
\end{minted}

Run this and you will see something like:

\begin{minted}{c}
i stores its value at 0x7ffeea3c4738
\end{minted}

The address is printed in hexadecimal.  Your address will differ, because
the operating system decides where in memory your variables live.

A \emph{pointer}\index{pointer} is a variable that holds an address.
If you want to store the address of an \texttt{int}, you declare a
variable of type \texttt{int *} (``pointer to int''):

\begin{minted}{c}
int i = 17;
int *addressOfI = &i;
printf("i is at %p\n", (void *)addressOfI);
\end{minted}

The \texttt{*} in the declaration is part of the type; it means ``this
variable holds an address, and the thing at that address is an
\texttt{int}.''

\subsection{Reading and writing through a pointer}

Once you have a pointer, you can use the \texttt{*} operator to read
or write the value \emph{at} the address the pointer holds.  This is
called \emph{dereferencing}\index{dereferencing} the pointer:

\begin{minted}{c}
int i = 17;
int *p = &i;

printf("%d\n", *p);   /* prints 17 -- reads through p  */

*p = 89;              /* writes 89 to the address p holds */
printf("%d\n", i);    /* prints 89 -- i was changed!   */
\end{minted}

To summarize:
\begin{itemize}
\item \texttt{\&i} gives you the address of \texttt{i}.
\item \texttt{int *p} declares a variable that can hold such an address.
\item \texttt{*p} gives you the value stored at the address \texttt{p} holds.
\end{itemize}

\subsection{NULL pointers}

Sometimes you want a pointer that does not yet point to anything.  Use
\texttt{NULL}\index{NULL}:

\begin{minted}{c}
int *p = NULL;   /* p points to nothing */

if (p) {
    printf("%d\n", *p);     /* only runs if p is non-null */
} else {
    printf("p is null\n");
}
\end{minted}

A null pointer is simply the address zero.  \textbf{Never dereference
a null pointer}, your program will crash.  Always check before
dereferencing a pointer that might be null.

\section{Pass-by-Reference}
\index{pass-by-reference}

Recall that when you call a function in C, arguments are \emph{copied}
into the function's frame.  But what if you want a function to modify a
variable in the calling function?  You pass the \emph{address} of that
variable instead of its value.  The function can then write to that
address.

There is a standard C function called \texttt{modf()}\index{modf} that
splits a floating-point number into its integer part and its fractional
part, and needs to return \emph{both}.  Since a function can only return
one value, it returns the fractional part and writes the integer part to
an address you supply:

\begin{minted}{c}
#include <stdio.h>
#include <math.h>

int main()
{
    double pi = 3.14;
    double integerPart;
    double fractionPart;

    /* Pass the address of integerPart so modf can write there */
    fractionPart = modf(pi, &integerPart);

    printf("Integer part:  %f\n", integerPart);
    printf("Fraction part: %f\n", fractionPart);
    return 0;
}
\end{minted}

Output:
\begin{minted}{c}
Integer part:  3.000000
Fraction part: 0.140000
\end{minted}

Here is another way to think about it.  Imagine you need to give a spy
an assignment.  You say: ``I've left a length of steel pipe at the foot
of the angel statue in the park.  When you have the information, roll
it up and leave it in the pipe.  I'll pick it up Tuesday.''  In the spy
business, this is called a \emph{dead drop}.

Pass-by-reference works the same way.  You are not giving the function
the data directly, you are giving it a \emph{location} where it can
leave the result.

Here is our own pass-by-reference function that converts meters to feet
and inches:

\begin{minted}{c}
#include <stdio.h>
#include <math.h>

void metersToFeetAndInches(double meters,
                            int *ftPtr,
                            double *inPtr)
{
    double rawFeet = meters * 3.281;
    int feet = (int)floor(rawFeet);
    double inches = (rawFeet - feet) * 12.0;

    if (ftPtr) *ftPtr = feet;
    if (inPtr) *inPtr = inches;
}

int main()
{
    int feet;
    double inches;
    metersToFeetAndInches(1.8, &feet, &inches);
    printf("%d ft %f in\n", feet, inches);
    return 0;
}
\end{minted}

The caller declares the variables it wants filled in (\texttt{feet},
\texttt{inches}), passes their addresses to the function, and the
function writes the results there.

\section{Structs: Grouping Related Data}
\index{struct}

A \emph{struct} lets you bundle several related values into a single
named type.  In Python you might use a dictionary for this; in C, a
struct is the fundamental grouping mechanism.

\begin{minted}{c}
struct Person {
    double heightInMeters;
    int weightInKilos;
};
\end{minted}

This defines a new type called \texttt{struct Person}.  You can then
create variables of that type:

\begin{minted}{c}
#include <stdio.h>

struct Person {
    double heightInMeters;
    int weightInKilos;
};

double bodyMassIndex(struct Person p)
{
    return p.weightInKilos / (p.heightInMeters * p.heightInMeters);
}

int main()
{
    struct Person mikey;
    mikey.heightInMeters = 1.70;
    mikey.weightInKilos  = 96;

    struct Person aaron;
    aaron.heightInMeters = 1.97;
    aaron.weightInKilos  = 84;

    printf("mikey BMI: %f\n", bodyMassIndex(mikey));
    printf("aaron BMI: %f\n", bodyMassIndex(aaron));
    return 0;
}
\end{minted}

You access the members of a struct using a dot: \texttt{mikey.heightInMeters}.

When \texttt{main} is running, the stack looks like this:

\begin{center}
\begin{tabular}{r|l|}
\hline
\textbf{main()} & \texttt{mikey.heightInMeters = 1.70} \\
                & \texttt{mikey.weightInKilos = 96} \\
                & \texttt{aaron.heightInMeters = 1.97} \\
                & \texttt{aaron.weightInKilos = 84} \\
\hline
\end{tabular}
\end{center}

The two \texttt{Person} structs live right inside \texttt{main}'s
frame on the stack.

\section{The Heap}
\index{heap}

The stack is fast and automatic, but it has two constraints: stack frames
are destroyed when functions return, and the sizes of all stack-allocated
variables must be known at compile time.

The \emph{heap} is a separate region of memory that has neither
constraint.  Memory on the heap persists until you explicitly release
it, and can be any size determined at runtime.

In C, you claim heap memory with \texttt{malloc}\index{malloc} and
release it with \texttt{free}\index{free}.

\begin{minted}{c}
#include <stdlib.h>

struct Person *mikey = malloc(sizeof(struct Person));
\end{minted}

\texttt{malloc} asks the system for enough heap memory to hold a
\texttt{Person} struct and returns the address of that memory.
The variable \texttt{mikey} is a \emph{pointer} to a \texttt{Person},
it lives on the stack, but the struct it points to lives on the heap.

\begin{figure}[h]
\centering
\begin{tabular}{|c|}
\hline
\textbf{The Stack} \\
\hline
\texttt{main()} \\
\texttt{mikey} = \framebox{\strut\quad} \\
\hline
\end{tabular}
\quad $\longrightarrow$ \quad
\begin{tabular}{|l|}
\hline
\textbf{The Heap} \\
\hline
\texttt{struct Person} \\
\texttt{heightInMeters = 1.70} \\
\texttt{weightInKilos = 96} \\
\hline
\end{tabular}
\caption{A pointer on the stack pointing to a struct on the heap}
\end{figure}

\subsection{The arrow operator}

When you have a pointer to a struct, you access its members using
\texttt{->}\index{arrow operator} instead of a dot:

\begin{minted}{c}
mikey->heightInMeters = 1.70;
mikey->weightInKilos  = 96;
\end{minted}

The code \texttt{mikey->weightInKilos} means: ``dereference the pointer
\texttt{mikey} to get the struct, then access its \texttt{weightInKilos}
member.''  It is shorthand for \texttt{(*mikey).weightInKilos}.

\subsection{Releasing heap memory}

When you are done with heap-allocated memory, you must release it with
\texttt{free}:

\begin{minted}{c}
free(mikey);
mikey = NULL;   /* good habit: prevents accidental reuse */
\end{minted}

After \texttt{free}, the memory is returned to the heap for reuse.
Setting the pointer to \texttt{NULL} prevents you from accidentally
using it again.

If you forget to call \texttt{free}, the memory is never returned.
This is called a \emph{memory leak}\index{memory leak}.  In a
short-lived program it may not matter, but in a long-running program
it can consume all available memory.

\subsection{A complete heap example}

\begin{minted}{c}
#include <stdio.h>
#include <stdlib.h>

struct Person {
    double heightInMeters;
    int weightInKilos;
};

double bodyMassIndex(struct Person *p)
{
    return p->weightInKilos / (p->heightInMeters * p->heightInMeters);
}

int main()
{
    struct Person *mikey = malloc(sizeof(struct Person));
    mikey->heightInMeters = 1.70;
    mikey->weightInKilos  = 96;

    printf("mikey BMI: %f\n", bodyMassIndex(mikey));

    free(mikey);
    mikey = NULL;

    return 0;
}
\end{minted}

\texttt{bodyMassIndex} now takes a \texttt{struct Person *} instead of
a \texttt{struct Person} by value.  This avoids copying the struct.
Functions that work on heap-allocated data almost always take pointers.

\subsection{Heap arrays}
\index{heap!arrays}

You can also allocate a contiguous block of memory for multiple values.
Use \texttt{malloc} with a count:

\begin{minted}{c}
/* Allocate space for 1000 doubles on the heap */
double *buffer = malloc(1000 * sizeof(double));

buffer[0] = 3.14;
buffer[1] = 2.72;

free(buffer);
buffer = NULL;
\end{minted}

A pointer to the first element is all you need to work with the whole
array.  The heap memory is contiguous, so \texttt{buffer[0]},
\texttt{buffer[1]}, and so on are laid out one after another.

\section{What Python Objects Actually Are}

You now know everything you need to understand what a Python object is,
under the hood.

When Python creates an object, it:
\begin{enumerate}
\item Allocates a struct on the heap large enough to hold the object's
  data (its instance variables).
\item Stores in that struct a pointer to the class's method table,
  so that method calls can find the right code.
\item Returns a pointer to that struct.
\end{enumerate}

Your Python variable \texttt{obj} is, at the machine level, a pointer
to a heap-allocated struct.  When Python's garbage collector determines
that nothing points to an object anymore, it calls the equivalent of
\texttt{free} to release that memory.

The next chapter introduces C++ classes, and with this memory model
in hand, you will implement linked lists, trees, and hash tables from
scratch.

\section{Exercises}

\begin{enumerate}

\item Write a C function \texttt{celsiusToFahrenheit} that takes a
  \texttt{double} temperature in Celsius and returns the equivalent in
  Fahrenheit.  Call it from \texttt{main} for the values 0, 20, 37,
  and 100, and print each result.

\item Write a function \texttt{swap} that takes two \texttt{int}
  pointers and swaps the values at the two addresses.  Test it so that
  after calling \texttt{swap(\&a, \&b)}, the values of \texttt{a} and
  \texttt{b} have been exchanged.

\item Write a struct \texttt{Point} with \texttt{double} members
  \texttt{x} and \texttt{y}.  Write a function \texttt{distance} that
  takes two \texttt{struct Point *} pointers and returns the Euclidean
  distance between them.  Allocate two points on the heap, fill in
  their coordinates, print the distance, and then free the memory.

\item \textbf{Memory leak hunt.}  The following program has a memory
  leak.  Identify and fix it.

\begin{minted}{c}
#include <stdio.h>
#include <stdlib.h>

struct Node {
    int value;
};

void printValue(int x)
{
    struct Node *n = malloc(sizeof(struct Node));
    n->value = x;
    printf("%d\n", n->value);
}

int main()
{
    for (int i = 0; i < 5; i++) {
        printValue(i);
    }
    return 0;
}
\end{minted}

\item A \texttt{struct tm} is defined in \texttt{<time.h>} and holds a
  broken-down calendar time.  The function \texttt{time(NULL)} returns
  the number of seconds since midnight, January 1, 1970.
  The function \texttt{localtime\_r(\&seconds, \&now)} fills a
  \texttt{struct tm} with the corresponding date and time.  Here is a
  snippet to get you started:

\begin{minted}{c}
#include <stdio.h>
#include <time.h>

int main()
{
    long secondsSince1970 = time(NULL);
    struct tm now;
    localtime_r(&secondsSince1970, &now);
    printf("The time is %d:%d:%d\n",
           now.tm_hour, now.tm_min, now.tm_sec);
    return 0;
}
\end{minted}

  Look up the other fields of \texttt{struct tm} and modify the
  program to print today's full date.  Then modify it to print the
  date that falls exactly 4,000,000 seconds from now.  (Hint:
  \texttt{tm\_mon} is 0-indexed, so January is 0.)

\end{enumerate}
